\begin{lemma}[Changes of frame]
Let $F$ be a finite-dimensional real inner-product space, let $0\le C_0<\infty$,
let $v\in\Aclass_{m,C_0}$ in the sense of \noderef{IsFitClass}, let
$z_0=(x_0,t_0)$ be a space-time point, let $R>0$ and let
$(p_0,a_0)\in F\times\R$. Let $\bar v$ be the change of frame of
\noderef{FrameChange},
\[
  \bar v(y,s)=
  \frac{v\bigl(x_0+R(y+p_0s),\,t_0+Rs\bigr)-v(x_0,t_0)}{R}
       -p_0\cdot y-\left(a_0+\frac{|p_0|^2}{2}\right)s .
\]
Then $\bar v\in\Aclass_{m,C_0}$. Moreover, writing $T(y,s)=(x_0+R(y+p_0s),
t_0+Rs)$ for the corresponding old point, if $(p,a)$ is a fit for $v$ at
$T(\zeta)$ of radius $R\rho$ with constant $C_0$, in the sense of
\noderef{IsFit}, then $(p-p_0,a-a_0)$ is a fit for $\bar v$ at $\zeta$ of radius
$\rho$ with constant $C_0$.
\end{lemma}
\begin{proof}
\emph{Step 1: the equation.} Write $\zeta=(y,s)$ for a space-time point in the
new frame, let
\[
  A(y,s)=-p_0\cdot y-\Bigl(a_0+\frac{|p_0|^2}2\Bigr)s
\]
be the affine part of the formula of \noderef{FrameChange}, and let $B$ be the
linear map $B(h,\ell)=(h+p_0\ell,\ell)$ on $F\times\R$. Then
$T(\zeta)=z_0+R\,B\zeta$, so the linear part of the affine map $T$ is $R\,B$, and
the formula of \noderef{FrameChange} reads
\begin{equation}\label{eq:fi-bar}
  \bar v=\frac{v\circ T-v(z_0)}R+A .
\end{equation}

\emph{(1a) $T$ is an affine homeomorphism, and $\bar v$ is continuous.}
The map $B$ is linear and invertible with $B^{-1}(h,\ell)=(h-p_0\ell,\ell)$,
because $B(h-p_0\ell,\ell)=(h,\ell)$ and $B^{-1}(h+p_0\ell,\ell)=(h,\ell)$; both
are continuous, being built from vector addition and scalar multiplication.
Since $R>0$,
the linear part $R\,B$ of $T$ is a linear homeomorphism of space-time onto
itself, so $T$ is an affine homeomorphism of space-time onto itself with
\[
  T^{-1}(Z)=B^{-1}\bigl(R^{-1}(Z-z_0)\bigr).
\]
In particular $T$ and $T^{-1}$ are continuous, $T$ maps open sets to open sets,
and $T$ carries every neighbourhood of a point $\zeta$ onto a neighbourhood of
$T(\zeta)$.

By clause (1) of \noderef{IsFitClass} the function $v$ is continuous on all of
space-time. In \eqref{eq:fi-bar} the composite $v\circ T$ is therefore
continuous, $v(z_0)$ is a constant, $R>0$, and $A$ is linear and hence
continuous; so $\bar v$ is continuous on all of space-time. The whole space is
open, which is the openness requirement that \noderef{IsLimitSubOn} and
\noderef{IsLimitSuperOn} place on $\Omega$ when $\Omega$ is the whole space.

\emph{(1b) The adjoint of $B$.}
The space-time inner product is
$\langle(h,\ell),(g,k)\rangle=\langle h,g\rangle+\ell k$, so for all $(h,\ell)$
and $(g,k)$
\[
  \bigl\langle B(h,\ell),(g,k)\bigr\rangle
   =\langle h+p_0\ell,g\rangle+\ell k
   =\langle h,g\rangle+\ell\bigl(\langle p_0,g\rangle+k\bigr)
   =\bigl\langle(h,\ell),\bigl(g,\ \langle p_0,g\rangle+k\bigr)\bigr\rangle ,
\]
whence
\begin{equation}\label{eq:fi-adjoint}
  B^{\mathsf T}(g,k)=\bigl(g,\ \langle p_0,g\rangle+k\bigr),
\end{equation}
which is the identity in the spatial slot. For a symmetric bilinear form $M$ on
space-time write $B^{\mathsf T}MB$ for the pulled-back form
$(\xi,\eta)\mapsto M[B\xi,B\eta]$; then, taking $\eta=\xi$,
\begin{equation}\label{eq:fi-quadform}
  \bigl(B^{\mathsf T}MB\bigr)[\xi]=M[B\xi]\qquad(\xi\in F\times\R).
\end{equation}

\emph{(1c) The frame change of a $C^2$ test, and the identity for $F$.}
Let $Z_0$ be a space-time point, put $\zeta_0=T^{-1}(Z_0)$, and let $\psi$ be
$C^2$ near $Z_0$. Define the frame change of $\psi$ by
\begin{equation}\label{eq:fi-phi}
  \phi=\frac{\psi\circ T-v(z_0)}R+A ,
\end{equation}
built from $\psi$ by the formula \eqref{eq:fi-bar}, with the additive constant
kept as $v(z_0)$. If $\psi$ is $C^2$ on
the open set $N\ni Z_0$, then $T^{-1}(N)$ is open and contains $\zeta_0$ by
(1a), and $\psi\circ T$ is $C^2$ there because $T$ is affine; since $A$ is
affine and $R>0$, the function $\phi$ is $C^2$ near $\zeta_0$.

Write $\nabla\psi(Z_0)=(\nabla_x\psi,\psi_t)$ for the space-time gradient of
$\psi$ at $Z_0$, split into spatial and time parts as in \noderef{LimitOp}, and
note that $A$ is linear with
$\nabla A=\bigl(-p_0,\,-(a_0+|p_0|^2/2)\bigr)$. Because $T$ is affine with
linear part $R\,B$, and the second derivative of an affine map vanishes, the
chain rule applied to \eqref{eq:fi-phi} gives, for all $\xi,\eta\in F\times\R$,
\[
  D\phi(\zeta_0)\xi=\frac{D\psi(Z_0)(R\,B\xi)}R+\langle\nabla A,\xi\rangle
   =D\psi(Z_0)(B\xi)+\langle\nabla A,\xi\rangle,
  \qquad
  D^2\phi(\zeta_0)[\xi,\eta]
   =\frac{D^2\psi(Z_0)[R\,B\xi,\,R\,B\eta]}R .
\]
The inner product identifies derivatives with gradients, so the first identity
reads $\langle\nabla\phi(\zeta_0),\xi\rangle
 =\langle\nabla\psi(Z_0),B\xi\rangle+\langle\nabla A,\xi\rangle
 =\langle B^{\mathsf T}\nabla\psi(Z_0)+\nabla A,\xi\rangle$ for every $\xi$, so
by \eqref{eq:fi-adjoint}
\[
  \nabla\phi(\zeta_0)=B^{\mathsf T}\nabla\psi(Z_0)+\nabla A
   =\Bigl(\nabla_x\psi-p_0,\ \langle p_0,\nabla_x\psi\rangle+\psi_t
      -a_0-\frac{|p_0|^2}2\Bigr) ,
\]
whose spatial part is
\begin{equation}\label{eq:fi-grad}
  \nabla_y\phi(\zeta_0)=\nabla_x\psi(Z_0)-p_0 .
\end{equation}
The second identity, since $R^{-1}R^2=R$, reads in the notation of
\eqref{eq:fi-quadform}
\begin{equation}\label{eq:fi-hess}
  D^2\phi(\zeta_0)=R\,B^{\mathsf T}\,D^2\psi(Z_0)\,B .
\end{equation}
These are the two identities that the change of frame induces, now asserted for
a $C^2$ test pair rather than for $v$ and $\bar v$, which by hypothesis are only
continuous, so that neither $D^2v$ nor $D^2\bar v$ need exist. Adding a constant to $\psi$ changes neither $\nabla\psi(Z_0)$ nor
$D^2\psi(Z_0)$, hence changes neither \eqref{eq:fi-grad} nor
\eqref{eq:fi-hess}: the particular constant $-v(z_0)/R$ in \eqref{eq:fi-phi} is
immaterial to what follows.

By \eqref{eq:fi-grad} and the definition of $B$,
\begin{equation}\label{eq:fi-Bxi}
  B\bigl(\nabla_y\phi(\zeta_0),1\bigr)
   =\bigl(\nabla_y\phi(\zeta_0)+p_0,\ 1\bigr)
   =\bigl(\nabla_x\psi(Z_0),1\bigr),
\end{equation}
so by the definition of $F$ in \noderef{LimitOp}, then \eqref{eq:fi-hess},
\eqref{eq:fi-quadform} and \eqref{eq:fi-Bxi},
\begin{equation}\label{eq:fi-Fscale}
  F[\phi](\zeta_0)=D^2\phi(\zeta_0)\bigl[(\nabla_y\phi(\zeta_0),1)\bigr]
   =R\,D^2\psi(Z_0)\bigl[B(\nabla_y\phi(\zeta_0),1)\bigr]
   =R\,D^2\psi(Z_0)\bigl[(\nabla_x\psi(Z_0),1)\bigr]
   =R\,F[\psi](Z_0).
\end{equation}

\emph{(1d) The test transfer, both signs.}
Let $\zeta_0$ be any space-time point and let $\phi$ be $C^2$ near $\zeta_0$
such that $\zeta\mapsto\bar v(\zeta)-\phi(\zeta)$ has a local maximum at
$\zeta_0$. Put $Z_0=T(\zeta_0)$ and
\begin{equation}\label{eq:fi-psi}
  \psi=R\,(\phi-A)\circ T^{-1}+v(z_0).
\end{equation}
If $\phi$ is $C^2$ on the open set $M\ni\zeta_0$, then $T(M)$ is open and
contains $Z_0$ by (1a), and $\psi$ is $C^2$ there because $T^{-1}$ and $A$ are
affine; so $\psi$ is $C^2$ near $Z_0$, and $T^{-1}(Z_0)=\zeta_0$.

First, $\phi$ is the frame change of this $\psi$ in the sense of
\eqref{eq:fi-phi}. Indeed, evaluating \eqref{eq:fi-psi} at $Z=T\zeta$, so that
$T^{-1}Z=\zeta$, gives $\psi(T\zeta)=R\bigl(\phi(\zeta)-A(\zeta)\bigr)+v(z_0)$,
whence
\[
  \frac{\psi(T\zeta)-v(z_0)}R+A(\zeta)
   =\bigl(\phi(\zeta)-A(\zeta)\bigr)+A(\zeta)=\phi(\zeta)
   \qquad(\zeta\in F\times\R).
\]
So (1c) applies to the pair $(\psi,\phi)$ at $Z_0$ and $\zeta_0$, and
\eqref{eq:fi-Fscale} holds for them.

Second, by \eqref{eq:fi-bar} evaluated at $\zeta=T^{-1}Z$ we have
$v(Z)=R\bigl(\bar v(T^{-1}Z)-A(T^{-1}Z)\bigr)+v(z_0)$, while
\eqref{eq:fi-psi} says $\psi(Z)=R\bigl(\phi(T^{-1}Z)-A(T^{-1}Z)\bigr)+v(z_0)$;
subtracting, the two copies of $R\,A(T^{-1}Z)$ and of $v(z_0)$ cancel and
\begin{equation}\label{eq:fi-diff}
  v(Z)-\psi(Z)=R\Bigl(\bar v\bigl(T^{-1}Z\bigr)-\phi\bigl(T^{-1}Z\bigr)\Bigr)
\end{equation}
for every space-time point $Z$.

Now transfer the local maximum. Let $N_0$ be a neighbourhood of $\zeta_0$ with
$\bar v-\phi\le(\bar v-\phi)(\zeta_0)$ on $N_0$. By (1a) the set $T(N_0)$ is a
neighbourhood of $Z_0$, and for $Z\in T(N_0)$ we have $T^{-1}Z\in N_0$, so by
\eqref{eq:fi-diff} and $R>0$
\[
  v(Z)-\psi(Z)=R\,\bigl(\bar v-\phi\bigr)\bigl(T^{-1}Z\bigr)
   \le R\,\bigl(\bar v-\phi\bigr)(\zeta_0)=v(Z_0)-\psi(Z_0),
\]
the last equality by \eqref{eq:fi-diff} at $Z=Z_0$. Hence
$Z\mapsto v(Z)-\psi(Z)$ has a local maximum at $Z_0$. By clause (1) of
\noderef{IsFitClass} and \noderef{IsLimitSolutionOn}, $v$ is a viscosity
subsolution of $F[v]=0$ on the whole space in the sense of
\noderef{IsLimitSubOn}; its subsolution inequality, applied at the point $Z_0$
with the test $\psi$, which is $C^2$ near $Z_0$ and for which $v-\psi$ has a
local maximum there, gives $F[\psi](Z_0)\ge0$. Since $R>0$,
\eqref{eq:fi-Fscale} yields
\[
  F[\phi](\zeta_0)=R\,F[\psi](Z_0)\ge0 .
\]

If instead $\bar v-\phi$ has a local \emph{minimum} at $\zeta_0$, the same
construction applies verbatim with the two inequalities reversed: $v-\psi$ has a
local minimum at $Z_0$ by \eqref{eq:fi-diff} and $R>0$; by clause (1) of
\noderef{IsFitClass} and \noderef{IsLimitSolutionOn}, $v$ is a viscosity
supersolution of $F[v]=0$ on the whole space in the sense of
\noderef{IsLimitSuperOn}, so its supersolution inequality at $Z_0$ with the test
$\psi$ gives $F[\psi](Z_0)\le0$; and \eqref{eq:fi-Fscale} yields
$F[\phi](\zeta_0)=R\,F[\psi](Z_0)\le0$.

\emph{(1e) $\bar v$ solves the equation.}
By (1a) the whole space is open and $\bar v$ is continuous on it. By (1d), for
every space-time point $\zeta_0$ and every $\phi$ that is $C^2$ near $\zeta_0$:
if $\bar v-\phi$ has a local maximum at $\zeta_0$ then $F[\phi](\zeta_0)\ge0$,
and if $\bar v-\phi$ has a local minimum at $\zeta_0$ then
$F[\phi](\zeta_0)\le0$. With $\Omega$ the whole space, a local maximum of
$\bar v-\phi$ on $\Omega$ at $\zeta_0$, as \noderef{IsLimitSubOn} requires, is
an ordinary local maximum at $\zeta_0$, and likewise for local minima in
\noderef{IsLimitSuperOn}. So the three conditions of each of those two
definitions hold, and $\bar v$ is a viscosity subsolution and a viscosity
supersolution of
$F[\bar v]=0$ on all of space-time, that is, by \noderef{IsLimitSolutionOn}, a
viscosity solution of $F[\bar v]=0$ there. Together with the continuity of
$\bar v$ this is clause (1) of \noderef{IsFitClass} for $\bar v$.

\emph{Step 2: the fits.}
Let $(p,a)$ be a fit for $v$ at $Z=T(\zeta)$ of radius $R\rho$. Let $h\in F$ and
$\ell\in\R$ with
\begin{equation}\label{eq:fi-newbox}
  |\ell|\le\rho,\qquad |h-(p-p_0)\ell|\le\rho .
\end{equation}
The increment $(h,\ell)$ in the new variables corresponds to the old increment
$R\,B(h,\ell)=\bigl(R(h+p_0\ell),R\ell\bigr)$, and
\[
  T(\zeta+(h,\ell))=Z+\bigl(R(h+p_0\ell),R\ell\bigr).
\]
The old box conditions for the fit $(p,a)$ at radius $R\rho$ are
$|R\ell|\le R\rho$ and $|R(h+p_0\ell)-pR\ell|\le R\rho$, i.e. $|\ell|\le\rho$
and $|h-(p-p_0)\ell|\le\rho$; these are exactly \eqref{eq:fi-newbox}. So
\noderef{IsFit} for $v$ gives
\begin{equation}\label{eq:fi-old}
  \Bigl|\,v\bigl(T(\zeta+(h,\ell))\bigr)-v(Z)
    -p\cdot R(h+p_0\ell)-\Bigl(a-\frac{|p|^2}2\Bigr)R\ell\,\Bigr|
  \le C_0R\rho .
\end{equation}
By the formula of \noderef{FrameChange},
\[
  \bar v(\zeta+(h,\ell))-\bar v(\zeta)
   =\frac{v\bigl(T(\zeta+(h,\ell))\bigr)-v(Z)}{R}
      -p_0\cdot h-\Bigl(a_0+\frac{|p_0|^2}2\Bigr)\ell .
\]
Subtracting the new affine increment
$(p-p_0)\cdot h+\bigl(a-a_0-|p-p_0|^2/2\bigr)\ell$ and using the identity
\[
  \Bigl(a_0+\frac{|p_0|^2}2\Bigr)+\Bigl(a-a_0-\frac{|p-p_0|^2}2\Bigr)
   =a-\frac{|p|^2}2+p\cdot p_0,
\]
which holds because
$|p-p_0|^2=|p|^2-2p\cdot p_0+|p_0|^2$, we obtain
\[
  \bar v(\zeta+(h,\ell))-\bar v(\zeta)-(p-p_0)\cdot h
     -\Bigl(a-a_0-\frac{|p-p_0|^2}2\Bigr)\ell
  =\frac1R\Bigl[\,v\bigl(T(\zeta+(h,\ell))\bigr)-v(Z)
     -p\cdot R(h+p_0\ell)-\Bigl(a-\frac{|p|^2}2\Bigr)R\ell\,\Bigr].
\]
Its absolute value is therefore at most $C_0R\rho/R=C_0\rho$ by
\eqref{eq:fi-old}. By \noderef{IsFit} this says that $(p-p_0,a-a_0)$ is a fit
for $\bar v$ at $\zeta$ of radius $\rho$ with constant $C_0$.

\emph{Step 3: $\bar v$ has fits at every center and radius.}
Let $\zeta$ be any space-time point and $\rho>0$. Since $v\in\Aclass_{m,C_0}$,
\noderef{IsFitClass} provides a fit $(p,a)$ for $v$ at $T(\zeta)$ of radius
$R\rho>0$. By Step 2, $(p-p_0,a-a_0)$ is a fit for $\bar v$ at $\zeta$ of radius
$\rho$. Together with Step 1 this shows $\bar v\in\Aclass_{m,C_0}$.
\end{proof}
